\documentclass[10pt,a4paper]{amsart}
\usepackage[margin=19mm]{geometry}
\usepackage{amsmath,amssymb,mathtools}
\usepackage[scr=rsfs,cal=euler]{mathalfa}
\usepackage{xfrac}
\usepackage[unicode,hidelinks,bookmarks=false]{hyperref}
\newcommand{\ie}{i.e.\ }
\newcommand{\cf}{cf.\ }
\newcommand{\resp}{respectively\ }
\newcommand{\Subsection}{\subsection}
\providecommand{\nobreakdash}{\nobreak\hbox{-}}
\setlength{\emergencystretch}{2em}
\multlinegap0pt
\usepackage{xcolor}
\usepackage{float}
\usepackage{amssymb}
\newtheorem{proposition}{Proposition}
\title{Topology of the Generalized Nash Equilibrium Problem}
\author{Matija Blagojevi\'c, Christof Sch\"utte}
\date{Source: arXiv:2507.03753v2 (CC BY 4.0)}
\begin{document}
\maketitle

\noindent\emph{Source and licence.} Topology of the Generalized Nash Equilibrium Problem. Version arXiv:2507.03753v2 (CC BY 4.0), distributed by arXiv under the Creative Commons Attribution 4.0 International licence, http://creativecommons.org/licenses/by/4.0/, as recorded in the arXiv licence metadata for this exact version and shown on the abstract page https://arxiv.org/abs/2507.03753v2. Full licence notice: \texttt{licenses/arxiv-2507.03753-CC-BY-4.0.txt}.\par\smallskip

\noindent\textit{Editorial scope.} Counted unit: the article's proposition prop:NI\_brings\_nothing\_new (source0.tex lines 438--441). The article's complete original proof of it is delivered with it (source0.tex lines 442--462, one \texttt{proof} environment of 21 lines). No mathematics is written, repaired or re-derived: the statement and the proof are the article's own lines, in the article's own notation, and no retained line is substituted or re-flowed.  No further source-local context is needed: every cross-reference of every retained line resolves inside the two delivered spans.  Nothing was omitted from any retained span, so the package carries no omission record.  No retained line contains a citation, so the wrapper carries no bibliography and no bibliography entry.  No retained line contains a figure environment, an image-inclusion command or drawing code, so no image is delivered and the package declares no asset dependency.  Editorial formatting only, in addition: an AMS \texttt{amsart} preamble that restates the article's own declarations and macros used by the retained lines, namely the environments \texttt{proposition} on separate counters, together with the standard packages those lines need.  The retained environments print in this document as Proposition 1 in order of appearance, whereas the article prints the same items with the numbering of its own full document.  The complete native source of the article as distributed in the arXiv e-print for this exact version is bundled unchanged as \texttt{sources/181041/source0.tex}.
\begin{proposition}\label{prop:NI_brings_nothing_new}
	Let $\mathcal{E}= (E_i,X_i,\theta_i : 1\leq i\leq N)$ be an abstract economy with $\theta_i\colon E\times E_i\longrightarrow\mathbb{R}$.
	Using the notation above, and denoting the global best-reply correspondence of $\mathcal{E}$ with $\Phi$ we have that $\Phi(x) = R(x)$ for every fixed point $x\in X(x)$.
\end{proposition}

\begin{proof}
	Instead of proving $\Phi(x)=R(x)$ we prove the equivalent statement that $X(x)\setminus\Phi(x)=X(x)\setminus R(x)$.

	\medskip\noindent
	First, assume that $y\in X(x)\setminus R(x)$.
	Consequently, by definition of $R(x)$ and $V(x)$, there exists $y'\in X(x)$ such that $\Psi(x,y') > \Psi(x,y)$, that is
	\[
		\sum_{i=1}^{N}\bigl(\theta_i(x,y'_i)-\theta_i(x,x_i)\bigr) > \sum_{i=1}^{N}\bigl(\theta_i(x,y_i)-\theta_i(x,x_i)\bigr) \quad \Longrightarrow \quad
		\sum_{i=1}^{N}\theta_i(x,y'_i) > \sum_{i=1}^{N}\theta_i(x,y_i).
	\]
	Hence, there exists at least one $1\leq j\leq N$ such that $\theta_j(x,y'_j) > \theta_j(x,y_j)$ and so $y_j\not\in \Phi_j(x)$ implying that  $y\in X(x)\setminus\Phi(x)$.

	\medskip
	Now assume that $y \in X(x)\setminus\Phi(x)$, or in other words there is an index $1\leq j\leq N$ such that $y_j\not\in \Phi_j(x)$.
	Consequently, there is $y_j'\in X_j(x)$ such that $\theta_j(x,y_j') > \theta_j(x,y_j)$ and so
	\begin{multline*}
		\Psi(x,y) = \Psi((x_1,\dots,x_j,\dots,x_N),(y_1,\dots,y_j,\dots,y_N)) = \sum_{i=1}^{N}\bigl(\theta_i(x,y_i)-\theta_i(x,x_i)\bigr) \\ =
		\sum_{i=1}^{j-1}\bigl(\theta_i(x,y_i)-\theta_i(x,x_i)\bigr) + (\theta_j(x,y_j)-\theta_j(x,x_j)) + \sum_{i=j+1}^{N}\bigl(\theta_i(x,y_i)-\theta_i(x,x_i)\bigr) \\ < \sum_{i=1}^{j-1}\bigl(\theta_i(x,y_i)-\theta_i(x,x_i)\bigr) + (\theta_j(x,y_j')-\theta_j(x,x_j)) + \sum_{i=j+1}^{N}\bigl(\theta_i(x,y_i)-\theta_i(x,x_i)\bigr)   \\ =  \Psi((x_1,\dots,x_j,\dots,x_N),(y_1,\dots,y_j',\dots,y_N)).
	\end{multline*}
	Therefore, $V(x)\neq\Psi(x,y)$ and so $x\in X(x)\setminus R(x)$.
\end{proof}

\end{document}
